% Avsnitt 15.1-15.6, ur lectures/L15/appendix/a_rx_shift_reg.md.

\section{Gränssnitt}
\label{c15:sec:iface}

\modulefig[0.62]{lectures/L15/appendix/images/rx_shift_reg.png}

\begin{booktable}{@{}p{6.8em}p{3.4em}p{10em}L@{}}
  \thead{Port} & \thead{Riktn.} & \thead{Typ} & \thead{Betydelse} \\
  \midrule
  \code{clock}, \code{reset_s2_n} & in & \code{std_logic} & 50 MHz klocka, synkroniserad aktiv låg
    reset. \\
  \code{sample} & in & \code{std_logic} & En puls per CAN-bit (\code{bit_timer.sample}) som markerar
    när \code{rx_bus} ska läsas. \\
  \code{rx_bus} & in & \code{std_logic} & Bussnivån, giltig vid sampelpunkten. \\
  \code{bit_count} & in & \code{std_logic_vector(3 downto 0)} & Antal \textbf{riktiga} bitar att
    samla in, 1--8; läses live, får ändras mellan grupper. \\
  \code{enable} & in & \code{std_logic} & Avgör om \code{sample}-pulser behandlas. \\
  \code{data} & out & \code{byte_t} & Den insamlade gruppen, \textbf{högerjusterad} i sina lägsta
    \code{bit_count} bitar. \\
  \code{valid} & out & \code{std_logic} & Pulsar när \code{bit_count} riktiga bitar samlats in. \\
  \code{stuff_error} & out & \code{std_logic} & Pulsar om en väntad stoppbit inte inverterar. \\
  \code{real_bit} & out & \code{std_logic} & Den enda avstoppade bit som accepterades vid det här
    samplet. \\
  \code{real_bit_valid} & out & \code{std_logic} & \code{'1'} vid ett sampel som accepterade en
    riktig bit; \code{'0'} vid en kastad stoppbit. \\
\end{booktable}

\code{data} och \code{real_bit} är nivåer; \code{valid}, \code{stuff_error} och
\code{real_bit_valid} är \textbf{pulser på en cykel}, satta till \code{'0'} som förval högst upp i
den klockade processen, precis som \code{tx_shift_reg}s. \code{real_bit_valid} är den som biter:
\code{can_controller} grindar \code{crc15}s enable direkt på den (\chapref{c17:ch}), så en nivå som
ligger kvar till nästa sampel matar in samma bit under en hel bitperiod och förstör CRC:n i tysthet.
Testbänken kontrollerar att den fallit igen cykeln efter varje sampel, av precis det skälet.

Det finns ingen \code{load}: samplingen sker mitt i perioden (\code{sample}, inte \code{bit_done}),
när alla noders signaler har hunnit stabilisera sig. En sändare ändrar bussen vid en periodgräns, så
en mottagare som också agerade på \code{bit_done} skulle läsa i samma ögonblick som nivån ändras;
\secref{c12:sub:bttrace}:s tickspår visar de två pulserna och de 14 ticken mellan dem.
\code{can_controller} ändrar bara \code{bit_count} mellan grupperna; \code{valid} rapporterar själv
varje grupps fullbordan. \code{rx_shift_reg_tb} binder \textbf{positionellt}, så ordningen och
typerna ovan är det som måste stämma; namnen är dina.

% ------------------------------------------------------------------------------------------
\section{Vad modulen måste minnas}
\label{c15:sec:state}

Som i \chapref{c14:ch}: lista tillståndet innan någon process skrivs. Mottagarsidan har samma två
uppgifter, speglade \textemdash{} hålla en grupp \emph{under påfyllning}, och hålla stoppbitsregeln
ärlig \textemdash{} plus en som den inte hade: att \emph{förutsäga} var sändaren måste ha satt in en
stoppbit. Varje regel nedan talar i termer av exakt de här fem signalerna:

\begin{vhdlcode}
signal data_reg      : byte_t;                     -- Collected real bits, filling from the right.
signal collected     : natural range 0 to 8;       -- Real bits accepted toward this chunk so far.
signal last_bit      : std_logic;                  -- The bit most recently sampled off the bus.
signal consecutive   : natural range 0 to MAX_RUN; -- Current run length; 0 = none.
signal stuff_expected: std_logic;                  -- The next sampled bit must be a stuff bit.
\end{vhdlcode}

\begin{itemize}
  \item \textbf{\code{data_reg} och \code{collected} är gruppen som fylls på.} Varje accepterad bit
    skiftas in från höger (\code{data_reg(6 downto 0) & rx_bus}), vilket är där högerjusteringen
    kommer ifrån, och \code{collected} räknar accepterade bitar mot \code{bit_count}. Registret bakom
    \code{data} måste vara den här interna signalen, med \code{data <= data_reg;} som en konkurrent
    tilldelning utanför processen: en \code{out}-port kan inte \emph{läsas}, och en skiftning är en
    läs-modifiera-skriv.
  \item \textbf{\code{last_bit} och \code{consecutive} är samma följdspårare som \chapref{c14:ch}:s},
    tecken för tecken: biten som senast låg på ledningen, och hur lång den identiska följd den
    avslutade är. Länkens båda sidor kör samma spårare över samma bitar, mot samma \code{MAX_RUN}
    från \code{can_def}; det är det som håller dem överens om var stoppbitarna hör hemma. Som i
    \chapref{c14:ch} skrivs paret från en enda procedur som processen deklarerar.
  \item \textbf{\code{stuff_expected} är spegelbilden av \chapref{c14:ch}:s \code{stuff_pending}}
    \textemdash{} samma flagga med motsatt uppgift. Sändarens flagga säger ”jag är \emph{skyldig} en
    stoppbit härnäst”; mottagarens säger ”nästa bit gör klokt i att \emph{vara} en”. Samma
    armeringsvillkor (en följd som når fem), samma tajming (kontrolleras vid nästa händelse), motsatt
    riktning på skyldigheten \textemdash{} och mottagarens version kan bli \emph{besviken}, och det
    är där \code{stuff_error} kommer ifrån.
\end{itemize}

Utgångarna registreras i samma process: \code{data} och \code{real_bit} är nivåer; \code{valid},
\code{stuff_error} och \code{real_bit_valid} är encykelspulserna som \secref{c15:sec:iface} beskrev.

% ------------------------------------------------------------------------------------------
\section{Processen, uppifrån och ned}
\label{c15:sec:process}

En klockad process håller allt ovanstående, plus en konkurrent tilldelning utanför den. Som i
\chapref{c14:ch} följer genomgången koden i den ordning den skrivs. \textbf{Följdspåraren} är rad för
rad \chapref{c14:ch}:s, med bara flaggan omdöpt efter vad den nu betyder: samma procedur, deklarerad
inuti den här modulens process och anropad från varje gren som tar en bit av ledningen utom en, och
undantaget noteras där det uppstår. Som i \chapref{c14:ch} jämförs \code{collected} mot
\code{bit_count} via \code{to_integer(unsigned(...))}, så filen behöver
\code{use ieee.numeric_std.all;}.

\textbf{Skelettet.} Den konkurrenta raden först: \code{data} är en \code{out}-port och kan inte
\emph{läsas}, medan en skiftning är en läs-modifiera-skriv, så registret är den interna
\code{data_reg} och porten speglar bara den. Sedan processen, med sin procedur deklarerad före
\code{begin}, asynkron reset först och allt annat på den stigande flanken:

\begin{vhdlcode}
data <= data_reg;  -- Concurrent, outside the process.

process(clock, reset_s2_n) is
    -- One procedure, the run tracker (below).
begin
    if (reset_s2_n = '0') then
        -- Clear everything.
    elsif (rising_edge(clock)) then
        -- Defaults, then the sample branch.
    end if;
end process;
\end{vhdlcode}

\subsection{Följdspåraren, skriven en gång}
\label{c15:sub:tracker}

\Chapref{c14:ch}:s procedur med \code{stuff_expected} i stället för \code{stuff_pending}; den enda
omdöpningen är hela skillnaden på mottagarsidan. Allt som \secref{c14:sub:tracker} säger om idiomet
(varför den läser signalerna direkt, varför den deklareras inuti processen i stället för i
arkitekturen, varför en \code{in}-parameter bevarar tajmingen före flanken, och varför syntesen
inlinar den) gäller här oförändrat:

\begin{vhdlcode}
    -- The run tracker: runs on the bit just taken off the wire.
    procedure track_run(new_bit: in std_logic) is
    begin
        if ((new_bit = last_bit) and (consecutive /= 0)) then
            consecutive <= consecutive + 1;
            if (consecutive = MAX_RUN - 1) then -- Run reaches MAX_RUN.
                stuff_expected <= '1';
            end if;
        else
            consecutive <= 1;                   -- Start a new run.
        end if;
        last_bit <= new_bit;
    end procedure;
\end{vhdlcode}

\subsection{Reset och förval}
\label{c15:sub:defaults}

\textbf{Reset nollställer allt}: gruppen (\code{data_reg}, \code{collected}), följdspåraren
(\code{last_bit}, \code{consecutive}, \code{stuff_expected}) och utgångarna. Det här är det
\textbf{enda} stället spåraren någonsin nollställs. Grupperna nollställer den inte: sändarens ström
är sammanhängande över hela ramen, så mottagarens bild av den måste vara det också.

\textbf{Förvalen härnäst, och bara de här tre}, så att grenarna nedan bara \emph{höjer} dem:

\begin{vhdlcode}
valid          <= '0';
stuff_error    <= '0';
real_bit_valid <= '0';
\end{vhdlcode}

\begin{warning}
\code{stuff_expected} hör \textbf{inte} hemma här, av exakt samma skäl som \code{stuff_pending} inte
hörde hemma i \chapref{c14:ch}:s förval: den är tillstånd, armerad vid det sampel som fullbordar en
följd och avgjord vid \emph{nästa} sampel, en hel bitperiod (femtio klockflanker) senare. Ge den
förvalet lågt så avdunstar förutsägelsen innan den kan förfalla, och då blir stoppbitar varken
kastade eller kontrollerade. Den skrivs på exakt tre ställen: nollställd vid reset, armerad inuti
\code{track_run}, och nollställd när förutsägelsen förfaller (båda utfallen nedan).
\end{warning}

\subsection{Samplingsgrenen}
\label{c15:sub:sample}

\textbf{\code{(sample = '1') and (enable = '1')} är den enda arbetande grenen}; det finns ingen
\code{load} på mottagarsidan. Den gör exakt en av tre saker, prövade i den här prioritetsordningen.
Först: en utestående förutsägelse förfaller, och den samplade biten avgör den åt det ena eller andra
hållet:

\begin{vhdlcode}
if ((sample = '1') and (enable = '1')) then
    if (stuff_expected = '1') then
        if (rx_bus /= last_bit) then
            stuff_expected <= '0';    -- The stuff bit arrived; prediction settled.

            track_run(rx_bus);        -- A stuff bit is a bit off the wire like any other.
        else
            stuff_error    <= '1';    -- Six identical bits: a stuffing violation.
            stuff_expected <= '0';
            consecutive    <= 1;      -- Restart the run at the violating bit; the
                                      -- tracker must never count past MAX_RUN.
        end if;
\end{vhdlcode}

Den anlända stoppbiten \emph{kastas}: ingen skiftning, ingen ändring av \code{collected}, och
\code{real_bit_valid} förblir \code{'0'}, eftersom en stoppbit aldrig är en riktig bit (det här är
biten som \code{crc15} inte får se). Anropet till spåraren är spegelbilden av \chapref{c14:ch}:s i
dess stoppgren: en kastad stoppbit resynkar spåraren precis som en utsänd gör, eftersom ledningen
inte skiljer riktiga bitar från stoppbitar och spåraren inte heller får göra det. Här behöver den
inget särskilt argument, eftersom \code{rx_bus /= last_bit} är den här grenens eget villkor, så
\code{track_run(rx_bus)} tar vägen för en ny följd per konstruktion och startar om räkningen på 1.

I \code{else}-grenen är sex identiska bitar i rad omöjliga på en korrekt stoppad buss, så
\code{stuff_error} pulsar. Det snyggaste för resten är att nollställa flaggan, inte acceptera något,
och starta om följdräkningen vid den brytande biten (\code{last_bit} håller redan dess värde),
eftersom det inte är den här modulens uppgift att avgöra vad ett brott \emph{betyder}. Omstarten är
inte valfri prydlighet: utan den ligger \code{consecutive} kvar på \code{MAX_RUN}, nästa identiska
bit kör spårarens \code{+ 1}, och signalens deklarerade intervall överskrids; testbänken matar in en
bit förbi brottet för att låsa fast precis det. Det här är också den enda placering
\code{track_run} inte kan tjäna, och av samma skäl: \code{rx_bus = last_bit} här, så anropet skulle
ta \code{+ 1}-vägen med \code{consecutive} redan på \code{MAX_RUN}. Den brytande biten är undantaget
från regeln proceduren kodar, så dess två rader står kvar utskrivna. \code{can_controller} behandlar
den som fatal och avbryter mottagningen (\chapref{c17:ch}), precis som en riktig kontroller skulle
göra.

I annat fall: acceptera en riktig bit, skifta in den från höger, exponera den bitvis för
\code{crc15}, avsluta gruppen om den här biten fullbordar den, och kör spåraren på biten som just
accepterats:

\begin{vhdlcode}
    else
        data_reg       <= data_reg(6 downto 0) & rx_bus;  -- Shift in from the right.
        real_bit       <= rx_bus;
        real_bit_valid <= '1';

        if (collected = to_integer(unsigned(bit_count)) - 1) then
            valid     <= '1';                 -- This bit completes the chunk.
            collected <= 0;
        else
            collected <= collected + 1;
        end if;

        track_run(rx_bus);   -- rx_bus was just accepted off the wire.
    end if;
end if;
\end{vhdlcode}

Jämförelsen här är läsningen av signalen före flanken från \chapref{c14:ch}: \code{collected} mot
\code{bit_count - 1} är idiomet ”följd-är-nu-full” tillämpat på gruppens fullbordan, samma form som
\code{consecutive = MAX_RUN - 1} inuti \code{track_run}. Ingenting annat händer vid en gruppgräns:
\code{bit_count} läses live, så nästa grupp börjar helt enkelt med nästa accepterade bit i den bredd
\code{can_controller} satt vid det laget, och \code{data_reg} nollställs \emph{inte}; den nya
gruppens bitar skiftas in under vad den gamla gruppen än lämnade ovanför dem
(Övning~\ref{c15:ex:justify} handlar om precis det).

Två spegelnoteringar värda att stanna vid:
\begin{itemize}
  \item \textbf{Det finns inget avslutande sampel.} \Chapref{c14:ch}:s \code{done} var tvungen att
    släpa en skiftning, eftersom en sänd bit ändå behövde sin fulla period på bussen. En mottagare
    samplar mitt i perioden, när biten redan fått sin stabiliseringstid, så i samma ögonblick som den
    sista riktiga biten accepteras är gruppen helt enkelt \emph{där}: \code{valid} går vid just det
    samplet, utan något att vänta på.
  \item \textbf{\code{data}/\code{valid} är gruppvisa, \code{real_bit}/\code{real_bit_valid}
    bitvisa.} En hel grupp är för grov för \code{crc15}, som behöver varje riktig bit när den
    anländer; \code{real_bit} och \code{real_bit_valid} exponerar exakt den enda accepterade biten,
    mottagarsidans motsvarighet till \code{tx_shift_reg}s par \code{bit_valid}/\code{stuff}.
    Motsvarighet, inte spegelbild: på sändarsidan bär stoppbitar fortfarande \code{bit_valid = '1'}
    och anroparen maskar bort dem med \code{stuff}, medan den kastade stoppbiten här aldrig höjer
    \code{real_bit_valid} alls, så \code{real_bit_valid} ensam är redan CRC-grindningen
    (\chapref{c17:ch}).
\end{itemize}

% ------------------------------------------------------------------------------------------
\section{En grupp, sampel för sampel}
\label{c15:sec:trace}

Att ta emot är inte helt enkelt \chapref{c14:ch} baklänges. En sändare \emph{bestämmer sig} för att
sätta in en stoppbit; en mottagare \emph{förutsäger} en och måste sedan kontrollera om den faktiskt
kom. Den skillnaden är där båda den här modulens intressanta utgångar kommer ifrån.

Här är ledningssekvensen som \code{tx_shift_reg} producerade för \code{"11111000"} i
\secref{c14:sec:trace}, anländande till ett nyss resettat \code{rx_shift_reg} med
\code{bit_count = 8}. Varje rad är en \code{sample}-puls. Kolumnen \code{följd} är antalet lika bitar
i rad efter samplet, och \code{data} visar registret fyllas nedifrån.

\begin{textdrawing}
            rx_bus     åtgärd  real_bit  real_bit_valid  följd      data  valid
sampel 1         1  acceptera         1               1     1  00000001      0
sampel 2         1  acceptera         1               1     2  00000011      0
sampel 3         1  acceptera         1               1     3  00000111      0
sampel 4         1  acceptera         1               1     4  00001111      0
sampel 5         1  acceptera         1               1     5  00011111      0
sampel 6         0      kasta         -               0     1  00011111      0
sampel 7         0  acceptera         0               1     2  00111110      0
sampel 8         0  acceptera         0               1     3  01111100      0
sampel 9         0  acceptera         0               1     4  11111000      1
\end{textdrawing}

Nio sampel för att återvinna åtta bitar, exakt de nio presentationer \chapref{c14:ch}:s sändare
behövde för att skicka dem. Tre rader bär lärdomen:
\begin{itemize}
  \item \textbf{Sampel 5} tar den femte \code{'1'}:an i rad. Ingenting kastas här. Att nå fem betyder
    bara att \emph{nästa} bit förväntas vara en stoppbit; den här biten är riktig och samlas in som
    vanligt.
  \item \textbf{Sampel 6} är den kastade stoppbiten. \code{real_bit_valid} faller till \code{'0'} och
    \code{data} rör sig inte: gruppen gör inga framsteg. Det här är raden som spelar roll för
    \chapref{c13:ch}:s motor, eftersom \code{crc15} inte får se den här biten. Följden resynkar till
    1 och tar stoppbitens eget värde som början på den nya följden.
  \item \textbf{Sampel 9} fullbordar den åttonde \emph{riktiga} biten, så \code{valid} pulsar och
    \code{data} läser \code{11111000}, den ursprungliga gruppen. Stoppbiten har försvunnit utan att
    \code{can_controller} någonsin fått veta att den fanns.
\end{itemize}

Lägg märke till att \code{data} fylls från höger, en position per accepterad bit. Efter nio sampel
råkar gruppen uppta alla åtta bitarna, men en 6-bitarsgrupp skulle sluta med sina bitar i
\code{data(5 downto 0)} och lämna det som ligger ovanför i fred. Det är högerjusteringen, och
Övning~\ref{c15:ex:justify} handlar om vad som sitter ovanför.

\textbf{När stoppbiten inte kommer.} Samma mekanism, en bit annorlunda. Sex \code{'0'}:or i rad där
den sjätte skulle ha inverterat:

\begin{textdrawing}
            rx_bus      väntat     åtgärd  real_bit_valid  följd  stuff_error
sampel 1         0  riktig bit  acceptera               1     1            0
sampel 2         0  riktig bit  acceptera               1     2            0
sampel 3         0  riktig bit  acceptera               1     3            0
sampel 4         0  riktig bit  acceptera               1     4            0
sampel 5         0  riktig bit  acceptera               1     5            0
sampel 6         0   stopp '1'      brott               0     1            1
\end{textdrawing}

Brottet upptäcks vid sampel \textbf{6}, inte sampel 5. Fem identiska bitar är fullt lagliga och krävs
till och med innan någon stoppbit alls existerar; det är den sjätte, som anländer där en inverterad
bit var skyldig, som är omöjlig på en korrekt stoppad buss. \code{stuff_error} pulsar under en cykel
och modulen fortsätter samla in; att avgöra vad ett brott \emph{betyder} hör till
\code{can_controller}, som avbryter ramen (\chapref{c17:ch}).

% ------------------------------------------------------------------------------------------
\section{Vad testbänken låser fast}
\label{c15:sec:tb}

\begin{itemize}
  \item Ledningssekvensen för \code{"11111000"} (stoppbit efter den femte \code{'1'}:an) avstoppas
    tillbaka till \code{"11111000"}, \code{valid} pulsar vid det sista samplet, och
    \code{real_bit_valid} är \code{'0'} vid den kastade stoppbiten och \code{'1'} (med
    \code{real_bit = rx_bus}) vid varje riktig bit.
  \item \code{real_bit_valid} är tillbaka på \code{'0'} cykeln \emph{efter} varje sampel, så en
    version som håller nivån underkänns även om den avstoppar perfekt.
  \item Sex \code{'1'}:or i rad höjer \code{stuff_error} vid det sjätte samplet, och en sjunde
    identisk bit accepteras som vanligt efteråt: följdräkningen startade om vid den brytande biten i
    stället för att räkna förbi \code{MAX_RUN}.
  \item En 6-bitarsgrupp pulsar \code{valid} vid sitt sjätte sampel, inte sitt åttonde, och landar
    högerjusterad i \code{data(5 downto 0)}.
  \item En 3-bitarsgrupp som matas in omedelbart därefter, \textbf{utan någon reset emellan}, plockar
    upp den nya \code{bit_count} och landar i \code{data(2 downto 0)}. Det här är den omdimensionering
    från grupp till grupp som \code{can_controller} utför mellan varje fält.
\end{itemize}

% ------------------------------------------------------------------------------------------
\section{Att lägga in den i \code{can_controller}}
\label{c15:sec:wire}

Det sista av de fyra CAN-blocken, och det första vars port map kopplas till något annat än klockan
och sina egna signaler:

\begin{vhdlcode}
    -- rx_shift_reg.
    signal rxsr_bit_count                     : std_logic_vector(3 downto 0);
    signal rxsr_enable                        : std_logic;
    signal rxsr_data                          : byte_t;
    signal rxsr_valid,    rxsr_stuff_error    : std_logic;
    signal rxsr_real_bit, rxsr_real_bit_valid : std_logic;
\end{vhdlcode}

\begin{vhdlcode}
    rx_shift_reg1: entity work.rx_shift_reg
        port map(clock, reset_s2_n, bt_sample, rx_bus_s2, rxsr_bit_count, rxsr_enable,
                 rxsr_data, rxsr_valid, rxsr_stuff_error, rxsr_real_bit, rxsr_real_bit_valid);
\end{vhdlcode}

Två av kopplingarna där är värda att stanna vid. \code{sample} kommer direkt från \code{bit_timer1}:s
utsignal, så mottagarvägen samplar på samma 70-procentstick som resten av designen tajmar mot;
ingenting däremellan bestämmer när bussen ska läsas. Och bussnivån är \code{rx_bus_s2}, utgången från
\chapref{c12:ch}:s instans \code{rx_bus_sync}, \textbf{inte} porten \code{rx_bus} själv.
\code{rx_shift_reg} är det enda block som läser bussen, och det läser den synkroniserade kopian;
ingenting i designen läser den råa pinnen. (Som med \code{tx_shift_reg}: ta
\code{reset_s2_n}-kopplingen som den är tills vidare; \chapref{c17:ch} flyttar båda skiftregistren
till en reset per ram.)

Med det här håller \code{can_controller}s arkitektur alla fyra CAN-delblocken och deras signaler, vid
sidan av \chapref{c12:ch}:s \code{meta_prev}, och inget rambeteende alls. Varje återstående koppling
mellan dem är \chapref{c16:ch}:s tillståndsmaskin.
